







\documentclass[11pt]{article}
\usepackage[T1]{fontenc}
\usepackage{lmodern}
\usepackage[margin=1in]{geometry}
\usepackage{amsmath,amssymb,amsthm}
\usepackage[hidelinks]{hyperref}
\newtheorem{rmk}{Remark}[section]
\usepackage{xcolor}

\numberwithin{equation}{section}
\newtheorem{theorem}{Theorem}[section]
\newtheorem{proposition}[theorem]{Proposition}
\newtheorem{definition}[theorem]{Definition}
\newtheorem{lemma}[theorem]{Lemma}


\newcommand{\br}[1]{\left( #1 \right)}

\renewcommand{\theequation}{%
	\ifnum\value{section}=0
	\arabic{equation}%
	\else
	\thesection.\arabic{equation}%
	\fi}

\renewcommand{\thetheorem}{%
	\ifnum\value{section}=0
	\arabic{theorem}%
	\else
	\thesection.\arabic{theorem}%
	\fi}

\newcommand{\fla}[1]{\textcolor{purple}{#1}}
\newcommand{\maxime}[1]{\textcolor{green!50!black}{#1}}


\title{Mode Stability for Charged Scalar Fields \\on  sub-extremal Reissner--Nordstr\"om black holes}
\author{Flavio Rossetti, Maxime Van de Moortel}
\date{\today}
\begin{document}
	\maketitle
	
	\begin{abstract}	
		We prove mode stability for massless charged scalar fields on 
		Reissner--Nordstr\"om black hole exterior regions in the full sub-extremal range, without  a smallness assumption on the charge coupling. This complements a previous result of Gajic on the extremal and near-extremal cases.\\
		
		\noindent\textbf{Formal verification and AI statement.}
		The proof was developed through interactions with AI models
		(ChatGPT 5.6 and Claude Opus 5.4), which were prompted to adapt
		the mode-stability approach introduced by P.~Hintz in
		\emph{Mode stability for the scalar wave equation on subextremal
			Kerr--de Sitter spacetimes} \cite{HintzKdS2026}. The argument presented here closely
		follows Hintz's approach.	The result answers a question of interest to the general relativity
		community. The key ideas underlying the adaptation were generated
		by the AI models; the authors have subsequently worked through
		the argument in detail and rewritten the proof, which will also be appended to our forthcoming proof of decay for charged scalar fields on  sub-extremal Reissner--Nordstr\"om black holes using Fourier-based methods \cite{RVdM}.	They take full responsibility for the correctness of the mathematical statements and
		proofs.		The present proof of mode stability has also been formally verified in Lean~4.
		The formalization is available \href{https://github.com/maxvdmmath/charged_scalar_field_mode_stability}{here}.
	\end{abstract}
	
	
	
	
	\section*{Introduction}
	
	We consider the massless charged wave equation with a fixed electromagnetic potential
	\begin{equation} \label{wave_eqn}
		g^{\mu\nu}D_\mu D_\nu\phi=0,\quad D_\mu=\nabla_\mu+iq_0A_\mu,
		\qquad A=\frac{e}{r}\,dt
	\end{equation}
	on the exterior of a sub-extremal Reissner--Nordstr\"om black hole,
	\begin{equation} \label{RN_metric}
		g=-\Omega^2dt^2+\Omega^{-2}dr^2+r^2d\omega_{\mathbb S^2}^2,
		\qquad
		\Omega^2=1-\frac{2M}{r}+\frac{e^2}{r^2}
		=\frac{(r-r_+)(r-r_-)}{r^2},
	\end{equation}
	where $M>0$, $0 < |e|<M$, $r_\pm=M\pm\sqrt{M^2-e^2}$, and $r>r_+$.
	Here $\nabla$ is the Levi--Civita connection of $g$, $\phi$ is complex-valued,
	and $q_0\in\mathbb R \setminus \{ 0 \}$ is its charge coupling. Both the metric $g$ and
	the electromagnetic potential $A$ are fixed. 
	
	The purpose of this note is to prove mode stability for every
	sub-extremal Reissner--Nordstr\"om exterior, every angular frequency, and every charge coupling. More precisely, we prove that any solution to \eqref{wave_eqn} of the form
	\begin{equation} \label{mode_solution}
		\phi(t, r, \theta, {\varphi}) = e^{-i\omega t} \frac{u(r)}{r}Y_L(\theta, {\varphi}),
	\end{equation}
	with $\Im \omega \ge 0$, $\omega \ne 0$,
	satisfying the standard outgoing regularity conditions at the event horizon and at infinity (see next paragraphs), has to be the trivial solution. Here $Y_L$ denotes {any nonzero spherical harmonic of degree $L$ on the sphere, satisfying}
	\[
	-\Delta_{\mathbb S^2}Y_L=L(L+1)Y_L,
	\qquad L\in\mathbb N_0.
	\]
	
	The proof is an adaptation
	of Hintz's argument for the scalar wave equation on
	Kerr--de Sitter~\cite{HintzKdS2026}. In that argument, two currents are constructed, one on each side of a frequency-dependent radius $r_0$ between the event horizon and the cosmological horizon. Each current has monotonicity properties and is associated to the solution conjugated by an oscillatory phase. Evaluated at $r_0$, the two currents give opposite-sign bounds on the same quantity, which forces the solution to vanish. This proves mode stability for real superradiant frequencies, which is the central part of Hintz's work.
	
	In the charged scalar field context considered here, the approach needs to be minimally modified to account for the asymptotically flat end. A key aspect, which does not appear in the cosmological setting, is a specific identity satisfied on the Reissner--Nordstr\"om black hole exterior, see already Lemma~\ref{lem:key-factorization}.	For sufficiently small $|q_0e|$, mode stability for a charged scalar field follows by perturbing
	the non-charged equation; see~\cite[Section~4]{Gajic2026I} and
	\cite[Section~1.1.7]{Gajic2026II}. Our result removes this smallness
	assumption throughout the sub-extremal range.
	
	To introduce the notation of the main statements, we consider the $(t, r, \theta, {\varphi})$ coordinates of \eqref{RN_metric}. We write $dx/dr=\Omega^{-2}$ and notice that the radial equation for $u$ (coming from \eqref{wave_eqn}) becomes
	\begin{equation} \label{radial_eqn}
		\partial_x^2u+
		\left[\left(\omega-\frac{q_0e}{r}\right)^2
		-\Omega^2\left(\frac{L(L+1)}{r^2}
		+\frac{\partial_r\Omega^2}{r}\right)\right]u=0.
	\end{equation}
	
	\begin{definition}[Mode solution] \label{def:mode_solution}
		We say that a {nonzero solution $\phi$ to \eqref{wave_eqn}} is a mode solution if 
		\[
		\phi(t, r, \theta, {\varphi}) = e^{-i\omega t} \frac{u(r)}{r}Y_L(\theta, {\varphi}),
		\]
		with $\omega \ne 0$,  where  $u(r)$ is such that
		\begin{align}
			u &= B_H e^{-i \br{\omega - \frac{q_0 e}{r_+}} x} \br{1 + o(1)}, \quad \text{ as } x \to -\infty, \label{bc_horizon} \\
			u &= B_I e^{i \omega x}r^{-i q_0 e}\br{1 + O( r^{-1})}, \quad \text{ as } x \to +\infty, \label{bc_infinity}
		\end{align}
		for some non-zero coefficients $B_H$ and $B_I$.
		We say that $\phi$ is an exponentially growing mode solution if the above is satisfied in the case $\Im(\omega) > 0$. We say that $\phi$ is a real mode solution if the above is satisfied with $\omega \in \mathbb{R} \setminus \{ 0\}$.
	\end{definition}
	
	
	The case $q_0e=0$ corresponds to the uncharged wave equation. We henceforth assume
	$q_0e>0$; the result for the opposite sign follows by sending
	$(\omega,q_0e)$ to $(-\overline\omega,-q_0e)$. Set
	\[
	\omega_{\mathrm{res}}:=\frac{q_0e}{r_+},
	\qquad \lambda:=q_0e,
	\qquad \sigma:=\frac{\omega}{\omega_{\mathrm{res}}},
	\qquad \kappa:=\frac{r_+-r_-}{r_+}.
	\]
	In the proof below, we divide $r$ and $x$ by $r_+$ and reuse the same
	notation for the rescaled coordinates. This gives the forthcoming 
	\eqref{eq:main-radial-equation-hintz}. 
	
	\begin{definition} Let  $0< \kappa \le 1$, $\lambda >0$ , $\Im \sigma \ge 0$, {$\sigma\ne0$}, $L \in \mathbb{N}_0$.	Let $u_H$ and $u_I$ be the unique\footnote{The existence and uniqueness statement is an application of results on ODE singularities, see e.g.\ \cite{PSRVdM_asymptotics, RVdM}.} solutions to \eqref{radial_eqn} satisfying, respectively, \eqref{bc_horizon} and \eqref{bc_infinity}{, with the normalizations $B_H=B_I=1$}. Then, we say that \emph{mode stability} holds for \eqref{wave_eqn} if the Wronskian 
		\[
		\mathcal W_{\kappa,\lambda}(\sigma,L) := \mathcal W(u_H, u_I) = \br{\partial_x u_H } u_I - u_H \br{ \partial_x u_I}
		\]
		satisfies
		\[
		\mathcal W_{\kappa,\lambda}(\sigma,L) \ne 0.
		\]
	\end{definition}
	Our theorem~\ref{thm:mode-stability-charged-scalar-field} states that mode stability holds for \eqref{wave_eqn}, i.e.\ that equation \eqref{wave_eqn}  does not admit any mode solution {with $\Im\omega\ge0$} in the sense of {Definition~\ref{def:mode_solution}}.
	
	
	
	The main obstruction to proving mode stability is \emph{superradiance}. In particular, suppose that $\phi$ is a real mode solution. Then, \eqref{radial_eqn} implies that  $\partial_x \Im \br{u' \overline{u}}=0$ (see already section \ref{section:real_not_superr}). Therefore, the boundary conditions \eqref{bc_horizon} and \eqref{bc_infinity} lead to the relation
	\begin{equation} \label{superradiant_condition}
		\omega |B_I|^2+
		(\omega-\omega_{\mathrm{res}})|B_H|^2=0.
	\end{equation}
	We define the superradiant interval
	\begin{equation} \label{superradiant_int}
		0<\omega<\omega_{\mathrm{res}}=\frac{q_0e}{r_+},
		\qquad\text{equivalently}\qquad 0<\sigma<1.
	\end{equation}
	Outside the superradiant interval, 
	\eqref{superradiant_condition} gives a contradiction, showing that no real mode solutions exist for such frequencies.
	In the superradiant interval \eqref{superradiant_int}, however, the two terms in \eqref{superradiant_condition}
	have opposite signs and no contradiction arises in such a range.
	Indeed, the flux of the conserved stationary energy into the horizon
	has the sign of $\omega(\omega-\omega_{\mathrm{res}})$. 
	The intuition behind this mechanism is that, although the metric is static,
	the interaction with the electric field allows the scalar field to
	extract energy from the black hole~\cite{Bekenstein1973}; see also
	the numerical study \cite{DiMenzaNicolas2015}. {Charged superradiance is related to the Klein phenomenon for electrostatic potentials, where amplification is compatible with an indefinite conserved energy; see Bachelot~\cite{Bachelot2004} for a rigorous scattering theory and a conjecture on mode stability for charged scalar fields.} In particular, a simple positivity argument does not suffice to prove mode stability for frequencies in the range \eqref{superradiant_int} for charged scalar fields, {unlike in the uncharged case}.
	
	Likewise, for
	$\Im\omega>0$, integrating $\partial_x \Im \br{u' \overline{u}}$ in $\mathbb{R}$ 
	confines a hypothetical growing mode to
	$0<\Re\omega<\omega_{\mathrm{res}}$; however it does not exclude it (see already lemma \ref{lem:localization-real-part}).
	
	
	The distinction between fixing and evolving the electromagnetic field
	is relevant here. In the Maxwell--charged-scalar-field system  the metric is fixed
	but $A$ evolves together with $\phi$ through the Maxwell equations,
	\[
	\nabla^\mu F_{\mu\nu}
	=\frac{iq_0}{2}
	\bigl(\phi\overline{D_\nu\phi}-\overline\phi D_\nu\phi\bigr),
	\qquad F=dA.
	\]
	This system is nonlinear, and the sum of the Maxwell and scalar
	stress-energy tensors is conserved. 
	In \cite{VdM2022}, boundedness, integrated local
	energy decay, and inverse-polynomial decay {were} proved for the above nonlinear system
	in spherical symmetry on fixed sub-extremal Reissner--Nordstr\"om
	and Schwarzschild exteriors, under smallness assumptions on the
	Maxwell charge and the initial scalar field energy.
	
	In the Einstein--Maxwell--charged-scalar-field system, the metric
	also evolves. The equation studied here is the scalar equation in the
	linearization of that system around Reissner--Nordstr\"om with
	$\phi=0$: the scalar stress-energy tensor and Maxwell current are
	quadratic in $\phi$, so scalar perturbations decouple at first order.
	The decay of charged fields in the exterior is also relevant to the
	interior problem and to Strong Cosmic Censorship. In spherical
	symmetry, the results of \cite{VdM2018}  {establish stability of the Cauchy horizon and curvature blow-up under quantitative decay assumptions on the event horizon}. The work \cite{KehleVdM2024}
	showed that, when this decay is slow, its oscillatory character can
	determine whether the metric and scalar field extend continuously
	across the Cauchy horizon. {On Reissner--Nordstr\"om--de Sitter, exponential local energy decay holds for massive charged Klein--Gordon fields when the charge product is sufficiently small at fixed positive mass~\cite{Besset2020}; for sufficiently small positive cosmological constant, exponentially growing modes occur in a regime of small charge and mass that includes the massless case~\cite{BessetHafner2021}.} Prescribing such a decay rate in the interior region allows  one to construct $H^1$ metric extensions beyond the Cauchy horizon of solutions to the spherically symmetric Einstein-Maxwell--charged scalar field system, when considering smooth initial data \cite{Rossetti2025}
	
	For the linear equation with fixed electromagnetic potential in the asymptotically flat case, Hod and
	Piran~\cite{HodPiran1998} predicted inverse-power tails, with
	oscillations along the event horizon. Gajic~\cite{Gajic2026I,Gajic2026II}
	recently proved energy boundedness, integrated energy decay, and
	precise late-time asymptotics on extremal and sufficiently
	near-extremal Reissner--Nordstr\"om black holes, without a smallness assumption
	on $|q_0e|$. Away from extremality, their extension to arbitrary
	coupling was conditional on real-frequency mode stability.
	Related late-time asymptotics for
	inverse-square potentials were obtained by Gajic~\cite{GajicInverseSquare2023}
	on Schwarzschild and by Gajic and the second author~\cite{GajicVdM2024}
	on Minkowski; the latter work also treats a charged wave equation
	with a fixed Coulomb field.
	
	Superradiance is also the principal obstruction to the elementary
	mode-stability argument for neutral waves on Kerr, where it is caused
	by rotation. Whiting~\cite{Whiting1989} famously overcame this obstruction
	for exponentially growing modes by an integral transformation of
	the radial equation, producing an equation to which a positive
	energy argument applies. Shlapentokh-Rothman~\cite{ShlapentokhRothman2015}
	proved real-axis mode stability and quantitative estimates using
	this transformation. {See also Andersson, Ma, Paganini, and Whiting~\cite{AnderssonMaPaganiniWhiting2017} for real-axis mode stability of the Teukolsky equation on sub-extremal Kerr.} Its adaptation to sub-extremal charged
	Reissner--Nordstr\"om does not remove superradiance completely:
	the transformed horizon frequency is a positive multiple of
	$\omega-q_0e/(r_++r_-)$, leaving
	$0<\omega<q_0e/(r_++r_-)$ untreated;
	see~\cite[Appendix~B]{Gajic2026I}.
	
	In the extremal case, Gajic~\cite[Theorem~7.14]{Gajic2026I}
	instead adapts the Whiting-type integral transformation introduced
	by Teixeira da Costa~\cite{TeixeiraDaCosta2020} for extremal Kerr.
	The resulting real-frequency Wronskian bounds, away from
	$\omega=0$ and $\omega=\omega_{\mathrm{res}}$, persist near
	extremality by an ODE perturbation argument. Separate estimates
	treat the threshold regimes. The remaining mode-stability problem
	at arbitrary coupling away from extremality is the one addressed
	here. On compact parameter and frequency sets away from the
	thresholds, the required positive lower bound for the Wronskian
	follows from nonvanishing and continuous dependence of the
	normalized outgoing solutions.
	
	Our proof follows Hintz's conjugation and current
	argument~\cite{HintzKdS2026}; the key steps are as follows: \begin{itemize}
		\item {Standard energy and flux identities reduce the analysis of nonzero-frequency modes with $\Im\sigma\ge0$ to $0<\Re\sigma<1$.}
		\item We compactify the black hole exterior, introducing the coordinate $y=1-\frac{r_+}{r} \in (0,1)$.
		\item  We remove the outgoing phases by setting $v_-=e^{i\Phi}u$ near
		the horizon $\{y=0\}$ and $v_+=e^{-i\Phi}u$ near infinity  $\{y=1\}$, with the two
		conjugations meeting at the transition point $y_0=1-\Re\sigma \in(0,1)$.
		\item  The conjugated $v_{\pm}$ admit currents which agree at $y=y_0$. For Reissner--Nordstr\"om, the sign
		needed to compare these currents follows from the explicit
		factorization in Lemma~\ref{lem:key-factorization}.
		\item The angular contribution becomes the constant $L(L+1)$; retaining
		it in the current gives additional nonnegative boundary terms.
		\item This proves mode stability for every $L$ without a restriction on
		$\lambda$. The same identity also excludes frequencies in the open
		upper half-plane.
		
		\item  The case $\sigma=0$ is treated (see 
		Proposition~\ref{prop:sigma-zero-threshold}) by  constructing  solutions with prescribed asymptotics near the regular singular point at infinity of the radial ODE; this is a separate argument, independent from the $\omega\neq 0 $ analysis.
	\end{itemize}
	
	
	
	
	
	\section{Mode stability for charged scalar fields}
	
	While the key argument in this section is strongly inspired by the 
	%two-sided phase
	%conjugation and current 
	construction in Hintz~\cite{HintzKdS2026}, the factorization used below is specific to the
	Reissner--Nordstr\"om coefficients considered here.
	
	\subsection{Reissner--Nordstr\"om coefficients and outgoing solutions}
	
	Fix $\lambda>0$ and 
	$L\in\mathbb N_0$, the latter corresponding to the  eigenvalue $L(L+1)$ for an $L$-th spherical harmonic.  We normalize the event horizon
	radius $$r_+=1,$$ and define $0<\kappa\leq1$ as twice the surface gravity of the event horizon of a Reissner-Nordstr{\"o}m black hole. Then, we set:
	\begin{equation}
		\Omega^2(r)
		:=\frac{(r-1)(r-1+\kappa)}{r^2},
		\qquad r\in(1,\infty),
		\label{eq:Omega2-definition}
	\end{equation}
	which coincides with  the function $\Omega^2$ in \eqref{RN_metric}, under our normalization assumptions.
	The tortoise coordinate $x$ is characterized, up to an additive
	constant, by
	\begin{equation}
		\frac{dr}{dx}=\Omega^2(r).
		\label{eq:tortoise-Omega}
	\end{equation}
	Thus $x\to-\infty$ at the event horizon $\{r=1\}$ and $x\to+\infty$ at infinity $\{r=+\infty\}$.
	We also set
	\begin{equation}
		K(r):=\frac{2-\kappa}{2r^3}-\frac{1-\kappa}{r^4}>0.
		\label{eq:K-and-g}
	\end{equation}
	
	In the following, we study
	\begin{equation}
		\partial_x^2u+
		\left[\lambda^2\left(\sigma-\frac1r\right)^2-2\Omega^2(r)K(r)
		-L(L+1)\frac{\Omega^2(r)}{r^2}\right]u=0.
		\label{eq:main-radial-equation-hintz}
	\end{equation}
	We note that \eqref{eq:main-radial-equation-hintz} coincides with \eqref{radial_eqn}, since the black hole mass and charge are related to $\kappa$ by the relations $2M = 2- \kappa$ and $e^2 = 1 - \kappa$.
	
	We then define the energy current $Q[u]:=\Im(\overline u\,\partial_xu)$: differentiating $Q[u]$ using \eqref{eq:main-radial-equation-hintz} then gives
	\begin{equation}
		Q[u]'(x)
		=-2\lambda^2\Im\sigma
		\left(\Re\sigma-\frac1r\right)|u|^2.
		\label{energy1}
	\end{equation}
	
	
	For $\Im\sigma\geq0$ we define the outgoing solution at the event horizon as the unique solution $u_H$ to \eqref{eq:main-radial-equation-hintz} satisfying
	\begin{equation}
		e^{i\lambda(\sigma-1)x}u_H=1+O(r-1),
		\qquad
		\partial_x\!\left(e^{i\lambda(\sigma-1)x}u_H\right)=O(r-1),
		\quad \text{ as }r\to1.
		\label{eq:horizon-outgoing-definition}
	\end{equation}
	
	
	For $\sigma\neq0$, the outgoing solution at infinity is defined using
	the long-range part of the potential in \eqref{eq:main-radial-equation-hintz}. Let $\Phi$ be defined by 
	\begin{equation}
		\partial_r\Phi
		= \frac{ \lambda(\sigma-r^{-1}) }{\Omega^2(r)}
		\label{eq:long-range-phase-derivative}
	\end{equation}
	Then 
	\begin{equation}
		\Phi(r)=\lambda\sigma r
		+\lambda\bigl((2-\kappa)\sigma-1\bigr)\log r+O(1)
		\text{ as  } r\to+\infty.
		\label{eq:long-range-phase-expansion}
	\end{equation}
	\begin{equation}
		\Phi(r)=\lambda(\sigma-1) x(r)\mathbin{{+}}{C_H}
		+O(e^{\kappa x(r)})
		\text{ as  } r\to1.
		\label{eq:short-range-phase-expansion}
	\end{equation}
	{Here $C_H\in\mathbb C$ depends on the additive constant in the phase.}
	The outgoing solution at infinity is defined as the unique solution $u_I$ to \eqref{eq:main-radial-equation-hintz} satisfying
	\begin{equation}
		e^{-i\Phi}u_I=1+O(r^{-1}),
		\qquad
		\partial_r\!\left(e^{-i\Phi}u_I\right)=O(r^{-2}), \quad \text{as } r \to+\infty.
		\label{eq:infinity-outgoing-definition}
	\end{equation}
	This coincides with \eqref{bc_infinity}, by exploiting the relation between the $r$ and $x$ coordinates in \eqref{eq:tortoise-Omega}.
	
	
	
	\begin{definition}[Outgoing at both ends]
		We say that a {nonzero} solution $u$ to \eqref{eq:main-radial-equation-hintz} is outgoing at both ends, {or, equivalently, corresponds to a mode solution}, if both \eqref{eq:horizon-outgoing-definition} and \eqref{eq:infinity-outgoing-definition} {hold with the leading constants $1$ replaced by independent nonzero amplitudes $B_H,B_I$}.
	\end{definition}
	
	Using the Wronskian convention
	\[
	W(f,h):=(\partial_x f) h-(\partial_x h)f,
	\]
	define
	\begin{equation}
		\mathcal W_{\kappa,\lambda}(\sigma,L):=W(u_H,u_I).
		\label{eq:Wronskian-hintz-section}
	\end{equation}
	The Wronskian is independent of $x$ (see \eqref{eq:main-radial-equation-hintz}) and its nonvanishing does not depend
	on the nonzero normalizations of $u_H$ and $u_I$.
	
	\begin{theorem}[Mode stability for the charged scalar field]
		\label{thm:mode-stability-charged-scalar-field}
		\label{thm:no-zeroes-upper-half-plane}
		For every $0<\kappa\leq1$, $\lambda>0$, and $L\in\mathbb N_0$,
		the radial equation \eqref{eq:main-radial-equation-hintz} has no
		nonzero solution outgoing at both ends for any $\sigma\neq0$ with
		$\Im\sigma\geq0$. Equivalently,
		\begin{equation}
			\mathcal W_{\kappa,\lambda}(\sigma,L)\neq0
			\qquad(\sigma\neq0,\ \Im\sigma\geq0).
			\label{eq:no-zeroes-nonzero-sigma}
		\end{equation}
		%At $\sigma=1$, the horizon condition selects the bounded solution.
		At $\sigma=0$, the same  conclusion holds with a specific convention for the outgoing solution at infinity\footnote{The family of outgoing solutions we construct is not continuous in the limit $\sigma \to 0^+$.} (see already  Proposition~\ref{prop:sigma-zero-threshold}).
		%This threshold value is defined separately.
	\end{theorem}
	
	\begin{rmk}
		During the preparation of this paper, we learned that  Z.~Li, S.~Ma, and L.~Zhang independently obtained the same mode stability result \cite{MaAl}. After a private communication between the two groups, we have jointly decided to make both papers public simultaneously.
	\end{rmk} 
	
	In the proof, we also use \eqref{eq:horizon-outgoing-definition} and
	\eqref{eq:infinity-outgoing-definition} with the leading constants $1$
	replaced by arbitrary amplitudes $B_H,B_I$, retaining the derivative
	bounds; at $\sigma=0$, we likewise allow $B_I=0$ in
	\eqref{eq:threshold-outgoing-definition}. The term outgoing continues
	to require a nonzero amplitude.
	
	
	\subsection{Localization of a hypothetical mode in the open upper half-plane}
	
	\begin{lemma}[Localization of the real part]
		\label{lem:localization-real-part}
		Suppose that $\Im\sigma>0$ and that a nonzero solution $u$ of
		\eqref{eq:main-radial-equation-hintz} is outgoing at both ends.
		Then $0<\Re\sigma<1$.
	\end{lemma}
	
	\begin{proof}
		The outgoing asymptotics \eqref{eq:horizon-outgoing-definition} and \eqref{eq:infinity-outgoing-definition} imply
		\begin{equation}
			\begin{aligned}
				|u|+|\partial_xu|&=O\!\left(e^{\lambda\Im\sigma\,x}\right)
				&& \text{as } x\to-\infty,\\
				|u|+|\partial_xu|&=O\!\left(e^{-\lambda\Im\sigma\,x}\right)
				&& \text{as } x\to+\infty,
			\end{aligned}
			\label{eq:outgoing-exponential-decay}
		\end{equation}
		where we used the fact that \eqref{eq:long-range-phase-derivative}  and \eqref{eq:tortoise-Omega} imply $\partial_x\Im\Phi=\lambda\Im\sigma$. 
		%Thus the long-range correction does not affect exponential decay in the $x$ coordinate.
		
		In particular, $Q[u]:=\Im(\overline u\,\partial_xu)$
		tends to zero  as $x\rightarrow \pm \infty$. Integrating \eqref{energy1} on $\mathbb{R}$ gives
		\begin{equation}
			2\lambda^2\Im\sigma
			\int_{-\infty}^{+\infty}\left(\Re\sigma-\frac1r\right)|u|^2\,dx=0
			\label{eq:imaginary-energy-identity}
		\end{equation}
		If $\Re\sigma\leq0$ or $\Re\sigma\geq1$, the integrand has a fixed
		sign, and since
		$0<r(x)^{-1}<1$, continuity therefore gives $u=0$, a contradiction.
		
	\end{proof}
	
	\subsection{Compactification of the exterior}
	
	Set
	\begin{equation}
		y:=1-\frac1r.
		\label{eq:y-compactification}
	\end{equation}
	Then $0<y<1$, with $y=0$ corresponding to the event horizon and $y=1$
	to infinity. Define
	\begin{equation}
		\mu(y):=\frac{dy}{dx}
		=\frac{\Omega^2(r)}{r^2}
		=y(1-y)^2\bigl[\kappa+(1-\kappa)y\bigr].
		\label{eq:mu-explicit}
	\end{equation}
	From now on, primes denote derivatives with respect to $y$.
	In particular,
	\begin{equation}
		\partial_x=\mu\partial_y,
		\qquad \partial_x^2u=\mu(\mu u')'.
		\label{eq:x-to-y-derivatives}
	\end{equation}
	Note that $\mu > 0$ for $0 < y <1$. Dividing \eqref{eq:main-radial-equation-hintz} by ${-}\mu$ gives
	\begin{equation}
		-(\mu u')'-\frac{\lambda^2(\sigma-1+y)^2}{\mu}u+[V+L(L+1)]u=0,
		\label{eq:y-equation}
	\end{equation}
	where (see \eqref{eq:K-and-g})
	\begin{align}
		V(y)&:=\frac{2\Omega^2(r(y))K(r(y))}{\mu(y)}=2r(y)^2K(r(y))\notag\\
		&=(2-\kappa)(1-y)-2(1-\kappa)(1-y)^2\notag\\
		&=(1-y)\bigl[\kappa+2(1-\kappa)y\bigr].
		\label{eq:V-explicit}
	\end{align}
	
	Thus
	\begin{equation}
		\mu(y)>0,\qquad V(y)>0\qquad0<y<1,
		\label{eq:V-positive}
	\end{equation}
	and
	\begin{equation}
		\mu(0)=\mu(1)=0,\qquad V(0)=\kappa,\qquad V(1)=0.
		\label{eq:V-endpoints}
	\end{equation}
	
	\subsection{Two-sided outgoing phase conjugation}\label{section:phase_conj}
	
	
	
	For the current argument, assume
	\begin{equation}
		\Im\sigma\geq0,\qquad 0<\Re\sigma<1,
		\qquad y_0:=1-\Re\sigma\in(0,1).
		\label{eq:current-parameter-range}
	\end{equation}  Note that $y_0$ corresponds to the location of the quadratic transition point in the potential, since when $\sigma \in \mathbb{R}$ \begin{equation}
		\sigma-\frac{1}{r(y_0)}=0.  
	\end{equation}
	Choose the additive constant in the phase in \eqref{eq:long-range-phase-derivative} so that
	\begin{equation}
		\Phi(y) =\lambda\int_{y_0}^{y}\frac{z+\sigma-1}{\mu(z)}\,dz.
		\label{eq:Phi-definition}
	\end{equation}
	Then
	\begin{equation}
		\Phi(y_0)=0,\qquad
		\mu(y)\Phi'(y)=\lambda(y-y_0)+i\lambda\Im\sigma, \quad \forall\, 0 < y < 1.
		\label{eq:phase-derivative}
	\end{equation}
	Note that the presence of a second-order transition point in \eqref{eq:y-equation}, due to the quadratic term $-\frac{\lambda^2(\sigma-1+y)^2}{\mu} u$, is the key difficulty. We address this difficulty by conjugating $u$ with its natural phase near the event horizon $\{y=0\}$ and near infinity $\{y=1\}$ respectively, namely we define
	\begin{equation}
		v_-(y):=e^{i\Phi(y)}u(y)\quad\text{for } 0<y\leq y_0,
		\qquad
		v_+(y):=e^{-i\Phi(y)}u(y)\quad\text{for }y_0\leq y<1.
		\label{eq:two-sided-conjugation}
	\end{equation}The corresponding phase defined as $\Phi(y) 1_{\{ 0<y\leq y_0 \}}- \Phi(y) 1_{\{ y_0 \leq y< 1 \}}$ is continuous at $y_0$ but not $C^2$; it is $C^1$ at $y_0$ if and only if $\Im(\sigma)=0$. 
	We define $y \mapsto \varepsilon(y)$ to be $\varepsilon=-1$ in $(0, y_0]$ and $\varepsilon=+1$ in $(y_0, 1)$,
	and, with a small abuse of notation, we write $u=e^{i\varepsilon\Phi}v_\varepsilon$ to denote either of the above functions. For a solution $u$ satisfying \eqref{eq:horizon-outgoing-definition} and \eqref{eq:infinity-outgoing-definition}, these conjugations
	remove the endpoint phases. 
	
	\begin{lemma}[Conjugated equation]
		\label{lem:conjugated-equation}
		The functions $v_\varepsilon$, $\varepsilon=\pm1$, satisfy
		\begin{equation}
			-(\mu v_\varepsilon')'+[V+L(L+1)]v_\varepsilon
			-i\varepsilon\lambda
			\bigl[2(y-y_0)v_\varepsilon'+v_\varepsilon\bigr]
			+2\varepsilon\lambda\Im\sigma\,v_\varepsilon'=0.
			\label{eq:conjugated-equation}
		\end{equation}
	\end{lemma}
	
	\begin{proof}
		Differentiating $u=e^{i\varepsilon\Phi}v_\varepsilon$ gives
		\begin{equation}
			\mu u'=e^{i\varepsilon\Phi}
			\bigl[\mu v_\varepsilon'
			+i\varepsilon\lambda(y+\sigma-1)v_\varepsilon\bigr].
			\label{eq:first-conjugated-derivative}
		\end{equation}
		\begin{align}
			(\mu u')'=e^{i\varepsilon\Phi}\biggl[
			(\mu v_\varepsilon')'
			+2i\varepsilon\lambda(y+\sigma-1)v_\varepsilon'
			+i\varepsilon\lambda v_\varepsilon\notag
			-\frac{\lambda^2(y+\sigma-1)^2}{\mu}v_\varepsilon\biggr].
			\label{eq:second-conjugated-derivative}
		\end{align}
		Substitute this into \eqref{eq:y-equation}. The quadratic 
		terms $\frac{\lambda^2(\sigma-1+y)^2}{\mu}u$ cancel, and writing
		$y+\sigma-1=(y-y_0)+i\Im\sigma$ gives
		\eqref{eq:conjugated-equation}.
	\end{proof}
	
	Combining $\partial_y=r^2\partial_r$ with  \eqref{eq:horizon-outgoing-definition} and \eqref{eq:infinity-outgoing-definition},  we find the following asymptotics: there are constants $B_H,B_I$
	such that
	\begin{equation}
		v_-(y)=B_H+O(y),\qquad v_-'(y)=O(1)
		\qquad\text{as }y\rightarrow 0,
		\label{eq:v-minus-endpoint}
	\end{equation}
	and
	\begin{equation}
		v_+(y)=B_I+O(1-y),\qquad v_+'(y)=O(1)
		\qquad\text{as }y\rightarrow 1.
		\label{eq:v-plus-endpoint}
	\end{equation}
	
	
	\subsection{The Reissner--Nordstr\"om positivity identity}
	
	Define
	\begin{equation}
		H(y):=2\mu(y)-(y-y_0)\mu'(y).
		\label{eq:H-definition}
	\end{equation}
	We will derive the following useful identity.
	\begin{lemma}
		\label{lem:key-factorization}
		For $0<\kappa\leq1$ and $0\leq y,y_0\leq1$,
		\begin{equation}
			H(y)-(y-y_0)^2V(y)
			=(1-y)(1-y_0)
			\bigl[\kappa(y+y_0-2yy_0)+2yy_0\bigr].
			\label{eq:key-factorization}
		\end{equation}
		In particular, if $0<y_0<1$, then
		\begin{equation}
			H(y)>(y-y_0)^2V(y)\geq0
			\qquad\text{for }0<y<1.
			\label{eq:H-strict-positive}
		\end{equation}
	\end{lemma}
	
	\begin{proof}
		{Expanding} \eqref{eq:mu-explicit},  \eqref{eq:V-explicit} and using the fact that
		\[
		\mu'(y) = (1-y)\br{ -4y^2(1-\kappa) - 3 \kappa y + 2y(1 - \kappa) + \kappa }
		\]
		{gives} \eqref{eq:key-factorization}.
		The right-hand side of \eqref{eq:key-factorization} is nonnegative for $0 < \kappa \le 1$ and $0 \le y, y_0 \le 1$. For
		$0<y,y_0<1$, the fact that
		\[
		y+y_0-2yy_0=y(1-y_0)+y_0(1-y)>0
		\]
		proves \eqref{eq:H-strict-positive}.
	\end{proof}
	
	\subsection{The modified Hintz current}
	
	Define, on each side of the transition point $y_0$, for $\varepsilon=\pm1$:
	\begin{align}
		\mathcal J_\varepsilon:={}&
		-\mu(y-y_0)|v_\varepsilon'|^2
		-\mu\Re\bigl(v_\varepsilon'\overline{v_\varepsilon}\bigr)
		\notag\\
		&+(y-y_0)[V+L(L+1)]|v_\varepsilon|^2
		+\varepsilon\lambda\Im\sigma\,|v_\varepsilon|^2.
		\label{eq:modified-Hintz-current}
	\end{align}
	
	\begin{lemma}[Differential current identity]
		\label{lem:differential-current-identity}
		For every $\varepsilon=\pm 1$, {on its corresponding half-interval}, we have:
		\begin{equation}
			\mathcal J_\varepsilon'
			=-\bigl[H+4\lambda\Im\sigma\,|y-y_0|\bigr]
			|v_\varepsilon'|^2
			+(y-y_0)V'|v_\varepsilon|^2.
			\label{eq:current-derivative}
		\end{equation}
	\end{lemma}
	
	\begin{proof}
		Multiply \eqref{eq:conjugated-equation} by
		$2(y-y_0)\overline{v_\varepsilon'}+\overline{v_\varepsilon}$
		and take real parts. 	The phase term contributes zero, since
		\[
		\Re\!\left(-i\varepsilon\lambda
		\left|2(y-y_0)v_\varepsilon'+v_\varepsilon\right|^2\right)=0.
		\]
		The remaining terms, after applying the product rule, give
		\[
		0=\mathcal J_\varepsilon'
		+\bigl[2\mu-(y-y_0)\mu'
		+4\varepsilon\lambda\Im\sigma\,(y-y_0)\bigr]
		|v_\varepsilon'|^2
		-(y-y_0)V'|v_\varepsilon|^2.
		\]
		Now use that $\varepsilon(y-y_0)=|y-y_0|$ on the entire interval $(0, 1)$.
	\end{proof}
	
	\begin{proposition}[Positive integrated identity]
		\label{prop:coercive-integrated-identity}
		Suppose $0<y_1<y_2<1$ and either $y_1 < y_2 \le y_0$ or $y_0 \le y_1 < y_2$.
		{Take $\varepsilon=-1$ if $y_2\le y_0$ and $\varepsilon=+1$ if $y_1\ge y_0$. Then:}
		\begin{align}
			&\mathcal J_\varepsilon(y_2)-\mathcal J_\varepsilon(y_1)
			-\left[(y-y_0)V|v_\varepsilon|^2\right]_{y_1}^{y_2}
			\notag\\
			&\quad=-\int_{y_1}^{y_2}
			\bigl[H-(y-y_0)^2V+4\lambda\Im\sigma\,|y-y_0|\bigr]
			|v_\varepsilon'|^2\,dy\notag\\
			&\qquad\ -\int_{y_1}^{y_2}
			V\left|v_\varepsilon+(y-y_0)v_\varepsilon'\right|^2\,dy.
			\label{eq:coercive-integrated-identity}
		\end{align}
		The RHS of \eqref{eq:coercive-integrated-identity} vanishes only if $v_\varepsilon$ vanishes on $[y_1,y_2]$.
		For outgoing solutions, the identity also holds at $y_1=0$ or
		$y_2=1$, with endpoint values understood as limits.
	\end{proposition}
	
	\begin{proof}
		We integrate by parts using $V$, keeping the constant angular
		term in $\mathcal J_\varepsilon$. This gives
		\begin{align}
			\int_{y_1}^{y_2}(y-y_0)V'|v_\varepsilon|^2\,dy
			={}\left[(y-y_0)V|v_\varepsilon|^2\right]_{y_1}^{y_2}
			-\int_{y_1}^{y_2}
			\left[V|v_\varepsilon|^2
			+2(y-y_0)V\Re\bigl(v_\varepsilon'
			\overline{v_\varepsilon}\bigr)\right]\,dy.
			\label{eq:V-integration-by-parts}
		\end{align}
		We then complete the square as follows:
		\begin{align}
			&\bigl[H+4\lambda\Im\sigma\,|y-y_0|\bigr]|v_\varepsilon'|^2
			+2(y-y_0)V\Re\bigl(v_\varepsilon'
			\overline{v_\varepsilon}\bigr)
			+V|v_\varepsilon|^2\notag\\
			&\quad=
			\bigl[H-(y-y_0)^2V+4\lambda\Im\sigma\,|y-y_0|\bigr]
			|v_\varepsilon'|^2
			+V\left|v_\varepsilon+(y-y_0)v_\varepsilon'\right|^2.
			\label{eq:complete-square}
		\end{align}
		Combining these two identities with \eqref{eq:current-derivative}
		proves \eqref{eq:coercive-integrated-identity}.
		
		By \eqref{eq:H-strict-positive}, \eqref{eq:V-positive}, 
		$\Im\sigma\geq0$ and $0< \Re \sigma <1$, if the sum of the
		integrals vanishes, continuity gives $v_\varepsilon'=0$ and
		$v_\varepsilon+(y-y_0)v_\varepsilon'=0$ throughout the interval;
		hence $v_\varepsilon=0$.
		
		Finally, the coefficients extend continuously to $y=0,1$, and
		\eqref{eq:v-minus-endpoint}--\eqref{eq:v-plus-endpoint} give bounded
		$v_\varepsilon$ and $v_\varepsilon'$ at the relevant endpoint.
		Thus the integrals converge and the identity extends to the
		endpoints by taking limits.
	\end{proof}
	
	\subsection{Evaluation of the current at endpoints \texorpdfstring{$y\in\{0,1\}$}{y = 0, 1} and the transition point \texorpdfstring{$y_0$}{y0}}
	
	At the transition point, $\Phi(y_0)=0$ by definition, so
	\begin{equation}
		v_\varepsilon(y_0)=u(y_0),\qquad
		v_\varepsilon'(y_0)
		=u'(y_0)+\varepsilon\frac{\lambda\Im\sigma}{\mu(y_0)}u(y_0).
		\label{eq:v-derivative-seam}
	\end{equation}
	Substitution into \eqref{eq:modified-Hintz-current} gives
	\begin{equation}
		\mathcal J_-(y_0)=\mathcal J_+(y_0)
		=-\mu(y_0)\Re\bigl(u'(y_0)\overline{u(y_0)}\bigr).
		\label{eq:seam-current-equality}
	\end{equation}
	Using \eqref{eq:v-minus-endpoint}--\eqref{eq:v-plus-endpoint} and
	\eqref{eq:V-endpoints}, we obtain
	\begin{equation}
		\begin{aligned}
			\mathcal J_-(0)
			&=-\bigl[y_0\bigl(\kappa+L(L+1)\bigr)
			+\lambda\Im\sigma\bigr]|B_H|^2,\\
			\mathcal J_+(1)
			&=\bigl[(1-y_0)L(L+1)+\lambda\Im\sigma\bigr]|B_I|^2.
		\end{aligned}
		\label{eq:endpoint-currents}
	\end{equation}
	
	We now apply \eqref{eq:coercive-integrated-identity} over the intervals $[0, y_0]$ (with $\varepsilon=-1$) and $[y_0, 1]$ (with $\varepsilon=1$). Let $\mathcal E_-$ and $\mathcal E_+$ denote the sums of the two
	nonnegative integrals in \eqref{eq:coercive-integrated-identity} over
	$[0,y_0]$ and $[y_0,1]$, respectively.
	In the case of the interval $[0, y_0]$, the boundary term is
	\[
	\left[(y-y_0)V|v_-|^2\right]_{0}^{y_0}
	=\kappa y_0|B_H|^2.
	\]
	It cancels the corresponding term in $\mathcal J_-(0)$, and hence
	\begin{equation}
		\mathcal J_-(y_0)
		=-\bigl[y_0L(L+1)+\lambda\Im\sigma\bigr]|B_H|^2-\mathcal E_-.
		\label{eq:left-seam-sign}
	\end{equation}
	When considering the interval $[y_0, 1]$ in \eqref{eq:coercive-integrated-identity}, the boundary term vanishes at both ends, so
	\begin{equation}
		\mathcal J_+(y_0)
		=\bigl[(1-y_0)L(L+1)+\lambda\Im\sigma\bigr]|B_I|^2+\mathcal E_+.
		\label{eq:right-seam-sign}
	\end{equation}
	Since the two currents coincide at $y_0$ (equation \eqref{eq:seam-current-equality}), we have:
	\begin{equation}
		\begin{aligned}
			0={}&\lambda\Im\sigma\bigl(|B_H|^2+|B_I|^2\bigr)\\
			&+L(L+1)\bigl[y_0|B_H|^2+(1-y_0)|B_I|^2\bigr]
			+\mathcal E_-+\mathcal E_+.
		\end{aligned}
		\label{eq:master-mode-identity}
	\end{equation}
	
	\subsection{Exclusion of exponentially growing modes and real superradiant modes}
	
	\begin{proposition}[No exponentially growing modes]
		\label{prop:no-upper-half-plane-modes}
		There is no nonzero solution of
		\eqref{eq:main-radial-equation-hintz} outgoing at both ends when
		$\Im\sigma>0$.
	\end{proposition}
	
	\begin{proof}
		If such a solution existed, Lemma~\ref{lem:localization-real-part}
		would give $0<\Re\sigma<1$, so $y_0 \in (0, 1)$ (see \eqref{eq:current-parameter-range}) and the results from section \ref{section:phase_conj} onwards apply.
		
		All terms in \eqref{eq:master-mode-identity} are nonnegative, so
		$\mathcal E_-=\mathcal E_+=0$. Proposition~\ref{prop:coercive-integrated-identity}
		then gives $v_-=v_+=0$, hence $u=0$, a contradiction.
		
		
	\end{proof}
	
	\begin{proposition}[No real modes in the superradiant interval]
		\label{prop:no-real-superradiant-modes}
		There is no nonzero solution of
		\eqref{eq:main-radial-equation-hintz} outgoing at both ends when
		$\sigma\in(0,1)$.
	\end{proposition}
	
	\begin{proof}
		Now $\Im\sigma=0$ and $0<y_0<1$. Equation
		\eqref{eq:master-mode-identity} gives
		\[
		L(L+1)\bigl[y_0|B_H|^2+(1-y_0)|B_I|^2\bigr]
		+\mathcal E_-+\mathcal E_+=0.
		\]
		All terms are nonnegative, so $\mathcal E_-=\mathcal E_+=0$. By the strict
		positivity in Proposition~\ref{prop:coercive-integrated-identity},
		$v_-=0$ on $(0,y_0)$ and $v_+=0$ on $(y_0,1)$. Thus $u=0$,
		a contradiction.
	\end{proof}
	
	\subsection{Real frequencies outside the superradiant interval}\label{section:real_not_superr}
	
	{For real $\sigma$, the coefficients in \eqref{eq:main-radial-equation-hintz} are real-valued.} Hence the energy current
	\begin{equation}
		Q[u]:=\Im\bigl(\overline u\,\partial_xu\bigr)
		\label{eq:usual-flux}
	\end{equation}
	is independent of $x$. If $\sigma\neq0$ and a solution
	satisfies both endpoint asymptotic conditions \eqref{eq:horizon-outgoing-definition} and \eqref{eq:infinity-outgoing-definition}, with amplitudes $B_H,B_I$
	possibly zero, then
	\begin{equation}
		Q[u](x)\to\lambda(1-\sigma)|B_H|^2
		\text{ as } x\to -\infty,
		\qquad
		Q[u](x)\mathrel{{\to}}\lambda\sigma|B_I|^2
		\text{ as } x\to +\infty.
		\label{eq:real-frequency-fluxes}
	\end{equation}
	
	
	For $\sigma<0$ or $\sigma>1$, this
	forces $B_H=B_I=0$. At $\sigma=1$, $u_H$ is
	bounded and satisfies
	\[
	\partial_xu_H\longrightarrow0\qquad x\to-\infty.
	\]
	Thus, $Q[u_H](x)\to 0$ as $x\to -\infty$, so equality with the
	$\lambda|B_I|^2$ term forces $B_I=0$. This also excludes
	a common outgoing solution. Together with
	Proposition~\ref{prop:no-real-superradiant-modes}, we conclude that
	\begin{equation}
		\mathcal W_{\kappa,\lambda}(\sigma,L)\neq0
		\qquad\text{for every real }\sigma\neq0.
		\label{eq:no-real-zeroes-except-threshold}
	\end{equation}
	
	\subsection{The case \texorpdfstring{$\sigma=0$}{sigma = 0}}
	
	At $\sigma=0$, the ODE \eqref{eq:main-radial-equation-hintz} becomes
	\begin{equation}
		\partial_x^2u+
		\left[\frac{\lambda^2-L(L+1)}{r^2}+O(r^{-3})\right]u=0 \text{ as }
		\qquad r\to+\infty.
		\label{eq:sigma-zero-asymptotic-equation}
	\end{equation}
	Equivalently, in the $r$ coordinate,
	\begin{equation}
		\partial_r^2u+O(r^{-2})\partial_ru
		+\left[\frac{\lambda^2-L(L+1)}{r^2}+O(r^{-3})\right]u=0.
		\label{eq:sigma-zero-r-equation}
	\end{equation}
	{The coefficients of the exact equation \eqref{eq:main-radial-equation-hintz} at $\sigma=0$, expressed in $r$, are analytic in $1/r$ near zero after the displayed powers are factored out.} 
	This ODE has a regular singular point at $r=+\infty$ (see e.g.\ \cite{Olver}), so there exists a sequence $(a_n)_{n \in \mathbb{N}} \subset {\mathbb{C}}$ so that
	\[
	u(r) = r^s \sum_{n=0}^{+\infty}a_n r^{-n}
	\]
	is a solution to {\eqref{eq:main-radial-equation-hintz} at $\sigma=0$} in $r > r_+$, where $s$ is one of the solutions to 
	\begin{equation}
		s(s-1)+\lambda^2-L(L+1)=0.
		\label{eq:threshold-indicial-equation}
	\end{equation}
	In the following, we make a choice for such an $s$, see already {Remark~\ref{rmk.choice}}.
	
	
	\begin{proposition}[Non-zero Wronskian when $\sigma=0$]
		\label{prop:sigma-zero-threshold}
		At $\sigma=0$, assume $B_I \ne 0$ and define the solution $u_I$ to {\eqref{eq:main-radial-equation-hintz}} as
		\begin{equation}
			u_I(r;\lambda,0)=
			\begin{cases}
				B_I r^{\frac12-\sqrt{(L+\frac12)^2-\lambda^2}}
				\bigl(1+O(r^{-1})\bigr),&0<\lambda<L+\frac12,\\[1mm]
				B_I r^{1/2}\bigl(1+O(r^{-1})\bigr),&\lambda=L+\frac12,\\[1mm]
				B_I r^{\frac12-i\sqrt{\lambda^2-(L+\frac12)^2}}
				\bigl(1+O(r^{-1})\bigr),&\lambda>L+\frac12.
			\end{cases}
			\label{eq:threshold-outgoing-definition}
		\end{equation}
		Here, the $O(r^{-1})$ term can be differentiated in $r$, and its $\partial_r$-derivative is  $O(r^{-2})$. With this definition, we have
		\begin{equation}
			\mathcal W_{\kappa,\lambda}(0,L)\neq0\qquad(\lambda>0).
			\label{eq:no-zero-at-sigma-zero}
		\end{equation}
	\end{proposition}
	
	\begin{proof}
		At $\sigma=0$, a nonzero solution  $u(x)$ satisfying the outgoing condition \eqref{eq:horizon-outgoing-definition} at the event horizon obeys the following asymptotics
		\begin{equation}
			u(x)=B_H e^{i\lambda x}\bigl(1+O(r(x)-1)\bigr), \text{ as }
			\ x\to-\infty,
			\label{eq:sigma-zero-horizon-asymptotic}
		\end{equation}
		with the corresponding differentiated asymptotics. Its energy current is
		therefore
		\begin{equation}
			Q[u]=\lambda|B_H|^2>0.
			\label{eq:sigma-zero-horizon-current}
		\end{equation}
		
		
		For $0<\lambda\leq L+\frac12$, write
		$u_I=B_Ir^s(1+\rho)$, where
		$s=\frac12-\sqrt{(L+\frac12)^2-\lambda^2}\leq\frac12$,
		$\rho=O(r^{-1})$ and $\partial_r\rho=O(r^{-2})$.
		Since $s$ is real,
		\[
		Q[u_I]
		=\Omega^2|B_I|^2r^{2s}
		\Im\bigl((1+\overline\rho)\partial_r\rho\bigr)
		=O(r^{2s-2})\longrightarrow0.
		\]
		By energy conservation, $Q[u_I]=0$, including when $B_I=0$.
		For $\lambda>L+\frac12$, the differentiable
		asymptotics in \eqref{eq:threshold-outgoing-definition}, together
		with $\partial_x=\Omega^2\partial_r$ and
		$\Omega^2\to1$, give
		\begin{equation}
			Q[u_I]
			=-\sqrt{\lambda^2-(L+\frac12)^2}\,|B_I|^2\leq0.
			\label{eq:sigma-zero-infinity-current}
		\end{equation}
		Thus a nonzero common outgoing solution would have positive flux
		at the horizon and nonpositive flux at infinity, contradicting
		conservation of $Q$. This proves
		\eqref{eq:no-zero-at-sigma-zero}.
	\end{proof}
	
	\begin{rmk}[Meaning of the threshold convention.]\label{rmk.choice}
		For $\lambda>L+\frac12$, {the powers}
		$r^{1/2\pm i\sqrt{\lambda^2-(L+1/2)^2}}$ have the same modulus. The sign in
		\eqref{eq:threshold-outgoing-definition} is therefore part of the
		definition, not a consequence of selecting the slower-growing
		solution. It selects the branch with negative flux; the selected
		branches for $0<\lambda\leq L+\frac12$ have zero flux.
		
		The value $\mathcal W_{\kappa,\lambda}(0,L)$ is defined separately by
		this convention.
	\end{rmk}
	
	
	\begin{proof}[Proof of Theorem~\ref{thm:mode-stability-charged-scalar-field}]
		Proposition~\ref{prop:no-upper-half-plane-modes} excludes all
		frequencies with $\Im\sigma>0$. The real superradiant interval
		is covered by Proposition~\ref{prop:no-real-superradiant-modes},
		and the flux argument excludes all other nonzero real frequencies.
		Proposition~\ref{prop:sigma-zero-threshold} treats $\sigma=0$
		under the stated threshold convention.
		
		
		More precisely, the preceding arguments give $B_I=0$ for $\sigma\neq0$
		and $B_H=0$ for $\sigma=0$ whenever a solution satisfies both
		asymptotic conditions \eqref{eq:horizon-outgoing-definition} and \eqref{eq:infinity-outgoing-definition}{, with \eqref{eq:threshold-outgoing-definition} replacing the latter when $\sigma=0$}.
		To deduce the Wronskian conclusion directly, suppose
		$W(u_H,u_I)=0$. Since $W(u_H,u_I)$ does not depend on $x$, it vanishes for every $x$.
		If $\sigma\neq0$, the  horizon outgoing condition \eqref{eq:horizon-outgoing-definition} makes $u_H$ nonvanishing
		sufficiently near the horizon, and
		\[
		\partial_x\!\left(\frac{u_I}{u_H}\right)
		=\frac{W(u_I,u_H)}{u_H^2}=0.
		\]
		Thus $u_I=cu_H$ there for some $c\in\mathbb C$. The solution $u_I$
		therefore satisfies the horizon outgoing condition \eqref{eq:horizon-outgoing-definition} with amplitude $cB_H$,
		while its  amplitude $B_I$ at $r=+\infty$ is nonzero, contradicting $B_I=0$.
		If $\sigma=0$, use instead that $u_I$ is nonvanishing sufficiently
		near infinity. The same quotient argument gives $u_H=cu_I$ there,
		so $u_H$ satisfies the threshold condition with amplitude $cB_I$.
		This contradicts $B_H=0$, since its horizon amplitude is nonzero.
		
		
	\end{proof}
	
	\begin{thebibliography}{99}
		
		\bibitem{HintzKdS2026}
		P.~Hintz,
		\emph{Mode stability for the scalar wave equation on subextremal
			Kerr--de Sitter spacetimes},
		arXiv:2608.25868v1 (2026).
		\url{https://arxiv.org/abs/2608.25868v1}.
		
		\bibitem{RVdM}
		F.~Rossetti and M.~Van de Moortel,
		\emph{Late-time asymptotics of charged scalar fields on sub-extremal Reissner--Nordstr{\"o}m black holes},
		in preparation.
		
		\bibitem{PSRVdM_asymptotics}
		F.~Pasqualotto, Y.~Shlapentokh-Rothman, and M.~Van de Moortel,
		\emph{The asymptotics of massive fields on stationary spherically symmetric black holes for all angular momenta},
		preprint (2023), \href{https://arxiv.org/abs/2303.17767}{arXiv:2303.17767} [gr-qc].
		
		\bibitem{Bekenstein1973}
		J.~D.~Bekenstein,
		\emph{Extraction of energy and charge from a black hole},
		Phys. Rev. D \textbf{7} (1973), 949--953.
		\url{https://doi.org/10.1103/PhysRevD.7.949}.
		
		\bibitem{Bachelot2004}
		{
			A.~Bachelot,
			\emph{Superradiance and scattering of the charged Klein--Gordon field by a step-like electrostatic potential},
			J. Math. Pures Appl. (9) \textbf{83} (2004), no.~10, 1179--1239.
			\url{https://doi.org/10.1016/j.matpur.2004.03.007}.}
		
		\bibitem{HodPiran1998}
		S.~Hod and T.~Piran,
		\emph{Late-time evolution of charged gravitational collapse and
			decay of charged scalar hair. II},
		Phys. Rev. D \textbf{58} (1998), 024018.
		
		\bibitem{DiMenzaNicolas2015}
		L.~Di Menza and J.-P.~Nicolas,
		\emph{Superradiance on the Reissner--Nordstr\"om metric},
		Class. Quantum Grav. \textbf{32} (2015), 145013.
		
		\bibitem{Besset2020}
		N.~Besset,
		\emph{Decay of the local energy for the charged Klein--Gordon equation in the exterior De Sitter--Reissner--Nordstr\"om spacetime},
		Ann. Henri Poincar\'e \textbf{21} (2020), no.~8, 2433--2484.
		\href{https://doi.org/10.1007/s00023-020-00919-z}{doi:10.1007/s00023-020-00919-z}.
		
		\bibitem{BessetHafner2021}
		{
			N.~Besset and D.~H\"afner,
			\emph{Existence of exponentially growing finite energy solutions for the charged Klein--Gordon equation on the de Sitter--Kerr--Newman metric},
			J. Hyperbolic Differ. Equ. \textbf{18} (2021), 293--310.
			\url{https://doi.org/10.1142/S0219891621500090}.}
		
		\bibitem{Rossetti2025}
		F.~Rossetti,
		\emph{Strong cosmic censorship for the spherically symmetric Einstein--Maxwell-charged-Klein--Gordon system with positive $\Lambda$: stability of the Cauchy horizon and $H^1$ extensions},
		Ann. Henri Poincar\'e \textbf{26} (2025), 675--753.
		\href{https://doi.org/10.1007/s00023-024-01454-x}{doi:10.1007/s00023-024-01454-x}.
		
		\bibitem{VdM2022}
		M.~Van de Moortel,
		\emph{Decay of weakly charged solutions for the spherically symmetric
			Maxwell-Charged-Scalar-Field equations on a Reissner--Nordstr\"om
			exterior space-time},
		Ann. Sci. \'Ec. Norm. Sup\'er. (4) \textbf{55} (2022), no.~2, 283--404.
		
		\bibitem{Gajic2026I}
		D.~Gajic,
		\emph{Charged scalar fields on Reissner--Nordstr\"om spacetimes I:
			integrated energy estimates},
		arXiv:2603.28874 (2026).
		\url{https://arxiv.org/abs/2603.28874}.
		
		\bibitem{Gajic2026II}
		D.~Gajic,
		\emph{Charged scalar fields on Reissner--Nordstr\"om spacetimes II:
			late-time tails and instabilities},
		arXiv:2603.28861 (2026).
		\url{https://arxiv.org/abs/2603.28861}.
		
		\bibitem{Olver}
		F.~W.~J. Olver,
		\emph{Asymptotics and Special Functions},
		AKP Classics,
		A K Peters, Wellesley, MA, 1997.
		Reprint of the 1974 original (Academic Press, New York).
		
		\bibitem{Whiting1989}
		B.~F.~Whiting,
		\emph{Mode stability of the Kerr black hole},
		J. Math. Phys. \textbf{30} (1989), 1301--1305.
		\url{https://doi.org/10.1063/1.528308}.
		
		\bibitem{ShlapentokhRothman2015}
		Y.~Shlapentokh-Rothman,
		\emph{Quantitative mode stability for the wave equation on the Kerr
			spacetime},
		Ann. Henri Poincar\'e \textbf{16} (2015), no.~1, 289--345.
		
		\bibitem{AnderssonMaPaganiniWhiting2017}
		{
			L.~Andersson, S.~Ma, C.~Paganini, and B.~F.~Whiting,
			\emph{Mode stability on the real axis},
			J. Math. Phys. \textbf{58} (2017), no.~7, 072501.
			\url{https://doi.org/10.1063/1.4991656}.}
		
		
		\bibitem{VdM2018}
		M.~Van de Moortel,
		\emph{Stability and instability of the sub-extremal
			Reissner--Nordstr\"om black hole interior for the
			Einstein--Maxwell--Klein--Gordon equations in spherical symmetry},
		Comm. Math. Phys. \textbf{360} (2018), 103--168.
		
		\bibitem{KehleVdM2024}
		C.~Kehle and M.~Van de Moortel,
		\emph{Strong Cosmic Censorship in the presence of matter:
			the decisive effect of horizon oscillations on the black hole
			interior geometry},
		Anal. PDE \textbf{17} (2024), no.~5, 1501--1592.
		
		\bibitem{GajicInverseSquare2023}
		D.~Gajic,
		\emph{Late-time asymptotics for geometric wave equations with
			inverse-square potentials},
		J. Funct. Anal. \textbf{285} (2023), no.~7, 110058.
		
		\bibitem{GajicVdM2024}
		D.~Gajic and M.~Van de Moortel,
		\emph{Late-time tails for scale-invariant wave equations with a
			potential and the near-horizon geometry of null infinity},
		arXiv:2401.13047 (2024).
		\url{https://arxiv.org/abs/2401.13047}.
		
		\bibitem{TeixeiraDaCosta2020}
		R.~Teixeira da Costa,
		\emph{Mode stability for the Teukolsky equation on extremal and
			subextremal Kerr spacetimes},
		Comm. Math. Phys. \textbf{378} (2020), no.~1, 705--781.
		
		\bibitem{MaAl}
		Z.~Li, S.~Ma, and L.~Zhang,
		\emph{Mode stability for the charged scalar field in
			Reissner--Nordstr\"om spacetime},
		in preparation
	\end{thebibliography}
\end{document}





